% Part: applied-modal-logics
% Chapter: epistemic-logic
% Section: public-announcement-logic-semantics
%READERNOTE{SOL6-A193: p ची घोषणा आणि p असूनही b ला p माहीत नाही याची घोषणा समान प्रतिमान देतात हा मूळ दावा चुकीचा होता. पहिल्यात दोन जगे, दुसऱ्यात फक्त w1 उरते; actual arrows व valuations वरून हे स्पष्ट केले. असत्य घोषणा-पूर्वांगात उरलेल्या प्रतिमानात मूल्यांकन करायचे नाही हेही स्पष्ट केले; म्हणून रिकामा restriction undefined point निर्माण करत नाही.}

\documentclass[../../../include/open-logic-section]{subfiles}

\begin{document}

\olfileid{aml}{el}{psm}

\olsection{सार्वजनिक घोषणा तर्कशास्त्राचे अर्थनिर्धारण}

सार्वजनिक घोषणा तर्कशास्त्रांची संबंधी प्रतिमाने
ज्ञानविषयक तर्कशास्त्रांतील प्रतिमानांसारखीच असतात.
मात्र सार्वजनिक घोषणा-संकारकाचे अर्थनिर्धारण नवे आहे.

\begin{defn}\ollabel{defn:mmodels} प्रतिमान~$\mModel M = \tuple{W, R, V}$
  मध्ये \emph{!!a{formula}~$!A$ हे~$w$ येथे सत्य असणे},
  म्हणजे $\mSat{M}{!A}[w]$, पुढीलप्रमाणे आगमनाने
  परिभाषित करतो:
  \begin{enumerate}
  \tagitem{prvFalse}{\indcase{!A}{\lfalse}{$\mSat{M}{\lfalse}[w]$ कधीच नसते}.}{}
  \tagitem{prvTrue}{\indcase{!A}{\ltrue}{$\mSat{M}{\ltrue}[w]$ नेहमी असते}.}{}
  \item $\mSat{M}{p}[w]$ तेव्हा आणि केवळ तेव्हाच, जेव्हा $w \in V(p)$
  \tagitem{prvNot}{\indcase{!A}{\lnot !B}{$\mSat{M}{\indfrm}[w]$
    तेव्हा आणि केवळ तेव्हाच, जेव्हा $\mSat/{M}{!B}[w]$}.}{}
  \tagitem{prvAnd}{\indcase{!A}{(!B \land !C)}{$\mSat{M}{\indfrm}[w]$
    तेव्हा आणि केवळ तेव्हाच, जेव्हा $\mSat{M}{!B}[w]$ आणि
    $\mSat{M}{!C}[w]$}.}{}
  \tagitem{prvOr}{\indcase{!A}{(!B \lor !C)}{$\mSat{M}{\indfrm}[w]$
    तेव्हा आणि केवळ तेव्हाच, जेव्हा $\mSat{M}{!B}[w]$ किंवा
    $\mSat{M}{!C}[w]$} (किंवा दोन्ही).}{}
  \tagitem{prvIf}{\indcase{!A}{(!B \lif !C)}{$\mSat{M}{\indfrm}[w]$
    तेव्हा आणि केवळ तेव्हाच, जेव्हा $\mSat/{M}{!B}[w]$ किंवा
    $\mSat{M}{!C}[w]$}.}{}
  \tagitem{prvIff}{\indcase{!A}{(!B \liff !C)}{$\mSat{M}{\indfrm}[w]$
    तेव्हा आणि केवळ तेव्हाच, जेव्हा $\mSat{M}{!B}[w]$ आणि
    $\mSat{M}{!C}[w]$ ही दोन्ही असतात, किंवा
    $\mSat{M}{!B}[w]$ आणि $\mSat{M}{!C}[w]$ यांपैकी
    कोणतेही नसते}.}{}
  \item\ollabel{defn:sub:mmodels-box}
    \indcase{!A}{\Knows_a !B}{$\mSat{M}{\indfrm}[w]$ तेव्हा आणि
    केवळ तेव्हाच, जेव्हा $R_a ww'$ असलेल्या प्रत्येक $w' \in W$
    साठी $\mSat{M}{!B}[w']$}
  \item\ollabel{defn:sub:mmodels-pal}
    \indcase{!A}{[!B] !C}{$\mSat{M}{\indfrm}[w]$ तेव्हा आणि
    केवळ तेव्हाच, जेव्हा एकतर $\mSat/{M}{!B}[w]$,
    किंवा $\mSat{M}{!B}[w]$ आणि $\mSat{M \mid !B}{!C}[w]$}


  येथे $\mModel M \mid !B = \tuple{W', R', V'}$ ची
  व्याख्या पुढीलप्रमाणे आहे:
  \begin{enumerate}
    \item $W' = \Setabs{ u \in W}{\mSat{M}{!B}[u]}$.
      म्हणजे $\mModel M \mid !B$ मधील जगे ही~$\mModel M$
      मध्ये $!B$ सत्य असलेली जगेच आहेत.
    \item $R'_a = R_a \cap (W' \times W')$. प्रत्येक कर्त्याचा
      प्राप्यता संबंध केवळ~$W'$ मध्ये उरलेल्या जगांपुरता
      मर्यादित केला जातो.
    \item $V'(p) = \Setabs{ u \in W'}{u \in V(p) }$.
      त्याचप्रमाणे उरलेल्या जगांतील विधानीय मूल्यांकने
      तीच राहतात. माहितीपर घटनांमुळे विधानीय चलांचे
      सत्यतामूल्य बदलत नाही, ही कल्पना त्यातून दिसते.
  \end{enumerate}
  $W'$ रिकामा असू शकतो. पण $!B$ हे~$w$ येथे असत्य
  असेल, तर घोषणा-सूत्र पहिल्या पर्यायाने सत्य आहे आणि
  त्या प्रतिबंधित रचनेत~$w$ येथे $!C$ तपासायचे नाही.
  $!B$ हे~$w$ येथे सत्य असेल, तर $w\in W'$ असल्याने
  $W'$ रिक्तेतर आहे आणि प्रतिबंधित प्रतिमानातील
  पुढील मूल्यांकन नीट परिभाषित आहे.

  \end{enumerate}
\end{defn}

म्हणून सार्वजनिक घोषणा तर्कशास्त्रांचे वैशिष्ट्य असे
की, प्रतिमान~$\mModel M$ येथे !!a{formula} सत्य आहे
की नाही हे कधीकधी~$\mModel M$ व्यतिरिक्त दुसरे
प्रतिमान पाहिल्याशिवाय ठरवता येत नाही.

हेही लक्षात घ्या की आपल्या अर्थनिर्धारणात घोषणा-संकारक
हा $\Box$~संकारक आहे. म्हणून एखाद्या जगात
!!a{formula}~$!A$ सत्य असल्याची घोषणा करता येत
नसेल, तर $[!A] !B$ तेथे आपोआप सत्य असते; जसे
अंतिम जगांत सर्व $\Box$ !!{formula}s सत्य असतात.

\begin{figure}
  \begin{center}
    \begin{tikzpicture}[modal]
      \node[world] (w1) [label={below left:$p, \lnot q$}]{$w_1$};
      \node[world] (w2) [left=of w1, label={left:$\lnot p, \lnot q$}]{$w_2$};
      \node[world] (w3) [above=of w1, label={right:$p, q$}] {$w_3$};
      \draw[<->] (w1) to [swap] node {$b$} (w2);
      \draw[->,loop below] (w1) to node {$a, b$} (w1);
      \draw[<->] (w1) to [swap] node {$a$} (w3);
      \draw[->,loop above] (w2) to node {$a, b$} (w2) ;
      \draw[->,loop above] (w3) to node {$a, b$} (w3) ;

      \node[world](v1)[right of=w1, xshift=1.25in, label={right:$p, \lnot q$}]{$w'_1$};
      \node[world](v3)[above of=v1,label={right:$p,q$}]{$w'_3$};
      \draw[<->] (v1) to node {$a$} (v3);
      \draw[->,loop below] (v1) to node {$a,b$} (v1);
      \draw[->,loop above] (v3) to node {$a,b$} (v3) ;

      \node(m)[below=of w1]{$\mModel M$};
      \node(m')[below=of v1]{$\mModel M \mid p$};

      \draw[-,dotted] (w1) to node {\footnotesize $p$ ची घोषणा}(v1) ;

    \end{tikzpicture}
  \end{center}
  \caption{$p$ ची सार्वजनिक घोषणा होण्यापूर्वी आणि झाल्यानंतर.}
  \ollabel{fig:announcement-example}
\end{figure}

!!a{formula} ची सार्वजनिक घोषणा म्हणजे प्रतिमान
संकुचित करणे, किंवा ज्या जगांत ते !!{formula} सत्य
होते तितक्याच जगांपुरते ते मर्यादित करणे, असे पाहता
येते. \olref{fig:announcement-example} मध्ये~$p$ ची
सार्वजनिक घोषणा केल्याचा परिणाम दाखवला आहे.
या प्रतिमानाचे एक वैशिष्ट्य असे की घोषणेमुळे
कर्ता~$b$ ला~$p$ माहीत होते, पण कर्ता~$a$ ला
तसे नव्याने कळत नाही (कारण~$p$ सत्य आहे हे
$a$ ला आधीच माहीत होते).

अधिक औपचारिकपणे, $\mSat{M}{\lnot \Knows_b p}[w_1]$
असते; पण $\mSat{M
\mid p}{\Knows_b p}[w'_1]$ असते. यावरून
$\mSat{M}{[p] \Knows_b
p}[w_1]$ मिळते. पण घोषणेच्या परिणामाबद्दल
याहूनही अधिक प्रबळ विधान करता येते. प्रत्यक्षात
$\mSat{M}{[p]\CKnows_{\{a,b\}} p}[w_1]$ असते.
म्हणजे $p$ ची घोषणा झाल्यावर ते \emph{सामायिक ज्ञान}
बनते.

पण हे सर्वसाधारणपणे खरे आहे का? $!A$ सत्य असल्याची
घोषणा केल्यावर $!A$ हे \emph{नेहमीच} सामायिक ज्ञान बनेल
का? आश्चर्य वाटेल, पण उत्तर नाही असे आहे. प्रत्यक्षात,
काही !!{formula}s सत्य असताना त्यांची घोषणा करणे
शक्य आहे, आणि घोषणेनंतर ती सत्य राहत नाहीत.
उदाहरणार्थ, \olref{fig:announcement-example} मधील
$w_1$ येथे $p \land \lnot \Knows_b p$ ची घोषणा
केल्याचा परिणाम पाहा. येथे $p$ हे $w_1,w_3$ या
दोन्ही जगांत सत्य आहे, म्हणून $\mModel M\mid p$ मध्ये
ही दोन्ही जगे उरतात. पण $w_3$ येथे $b$ चा एकमेव
प्राप्य पर्याय $w_3$ असल्याने $\Knows_b p$ सत्य आहे;
त्यामुळे $p\land\lnot\Knows_b p$ तेथे असत्य आहे.
हे संपूर्ण सूत्र फक्त $w_1$ येथे सत्य असते. म्हणून
$\mModel M\mid(p\land\lnot\Knows_b p)$ मध्ये केवळ
$w_1$ उरते आणि दोन्ही कर्त्यांसाठी फक्त त्याचे स्वबाण
उरतात; ते $\mModel M\mid p$ सारखे प्रतिमान नाही.
या एक-जगी प्रतिमानात $p$ आणि $\Knows_b p$ दोन्ही
सत्य आहेत. म्हणून
$\mSat{M \mid (p \land \lnot \Knows_b p)}{\lnot
(p \land \lnot \Knows_b p)}[w_1]$ असते. म्हणजे
घोषणा केल्यावर असत्य होणारे हे !!a{formula} आहे.

\end{document}
